\documentclass[
	% -- opções da classe memoir --
	12pt,				% tamanho da fonte
	openright,			% capítulos começam em pág ímpar (insere página vazia caso preciso)
	oneside,			% twoside para impressão em recto e verso. Oposto a oneside
	a4paper,			% tamanho do papel. 
	% -- opções da classe abntex2 --
	chapter=TITLE,		% títulos de capítulos convertidos em letras maiúsculas
	section=TITLE,		% títulos de seções convertidos em letras maiúsculas
	%subsection=TITLE,	% títulos de subseções convertidos em letras maiúsculas
	%subsubsection=TITLE,% títulos de subsubseções convertidos em letras maiúsculas
	% -- opções do pacote babel --
	english,			% idioma adicional para hifenização
	french,				% idioma adicional para hifenização
	spanish,			% idioma adicional para hifenização
	brazil,				% o último idioma é o principal do documento
 ]{abntex2}

% ---
% PACOTES
% ---

% ---
% Pacotes fundamentais
% ---
\usepackage{lmodern}
\usepackage[T1]{fontenc}
\usepackage[utf8]{inputenc}
\usepackage{indentfirst}
\usepackage{color}
\usepackage{graphicx}
\usepackage{microtype}
\usepackage{tabularx}

% ---
% Pacotes adicionais
% ---
\usepackage{multicol}
\usepackage{multirow}
\usepackage{lipsum}				% para geração de dummy text
% ---

% ---
% Pacotes de citações
% ---
\usepackage[brazilian,hyperpageref]{backref}	 % Paginas com as citações na bibl
\usepackage[alf]{abntex2cite}	% Citações padrão ABNT

% --- 
% CONFIGURAÇÕES DE PACOTES
% --- 

% ---
% Configurações do pacote backref
% Usado sem a opção hyperpageref de backref
\renewcommand{\backrefpagesname}{Citado na(s) página(s):~}
% Texto padrão antes do número das páginas
\renewcommand{\backref}{}
% Define os textos da citação
\renewcommand*{\backrefalt}[4]{
	\ifcase #1 %
		Nenhuma citação no texto.%
	\or
		Citado na página #2.%
	\else
		Citado #1 vezes nas páginas #2.%
	\fi}%
% ---

% ---
% Informações de dados para CAPA e FOLHA DE ROSTO
% ---
\titulo{Segurança Alimentar e Agronegócio Sustentável no Estado de São Paulo}
\autor{Cícero Santos Melo \\ Ivan Lacava Júnior \\ Góis Polícia Civil \\ Márcio Bontempo \\ Marcos Roberto Silva \\ Pedro Octavio de Souza Barros}
\local{São Paulo}
\data{2026}
\instituicao{%
  ASSOCIAÇÃO DOS DIPLOMADOS DA \\ ESCOLA SUPERIOR DE GUERRA
  \par
  Delegacia do Estado de São Paulo
  }
\tipotrabalho{Trabalho de Conclusão de Curso}
% O preambulo deve conter o tipo do trabalho, o objetivo, 
% o nome da instituição e a área de concentração 
\preambulo{Trabalho de Conclusão de Curso apresentado à
Associação dos Diplomados da Escola Superior de
Guerra, Delegacia do Estado de São Paulo, como
requisito para a obtenção do diploma do Curso de
Estudos de Política e Estratégia, 69a turma.}
\orientador{Prof. Reinaldo Bento de Lima}
% ---

% ---
% Configurações de aparência do PDF final

% alterando o aspecto da cor azul
\definecolor{blue}{RGB}{41,5,195}

% informações do PDF
\makeatletter
\hypersetup{
     	%pagebackref=true,
		pdftitle={\@title}, 
		pdfauthor={\@author},
    	pdfsubject={\imprimirpreambulo},
	    pdfcreator={LaTeX with abnTeX2},
		pdfkeywords={abnt}{latex}{abntex}{abntex2}{relatório técnico}, 
		colorlinks=true,       		% false: boxed links; true: colored links
    	linkcolor=blue,          	% color of internal links
    	citecolor=blue,        		% color of links to bibliography
    	filecolor=magenta,      		% color of file links
		urlcolor=blue,
		bookmarksdepth=4
}
\makeatother
% --- 

% --- 
% Espaçamentos entre linhas e parágrafos 
% --- 

% O tamanho do parágrafo é dado por:
\setlength{\parindent}{1.3cm}

% Controle do espaçamento entre um parágrafo e outro:
\setlength{\parskip}{0.2cm}  % tente também \onelineskip

% ---
% compila o indice
% ---
\makeindex
% ---

% ----
% Início do documento
% ----
\begin{document}

% Seleciona o idioma do documento (conforme pacotes do babel)
%\selectlanguage{english}
\selectlanguage{brazil}

% Retira espaço extra obsoleto entre as frases.
\frenchspacing 

% ----------------------------------------------------------
% ELEMENTOS PRÉ-TEXTUAIS
% ----------------------------------------------------------
% \pretextual

% ---
% Capa
% ---
\imprimircapa
% ---

% ---
% Folha de rosto
% (o * indica que haverá a ficha bibliográfica)
% ---
\imprimirfolhaderosto*
% ---

% ---
% Ficha catalográfica
% Substitua os campos provisórios pelos dados fornecidos pela biblioteca.
% ---
\begin{fichacatalografica}
\begin{center}
\fbox{\begin{minipage}{12cm}
\small
\textbf{Ficha catalográfica}

\vspace{0.5cm}

SANTOS MELO, Cícero; LACAVA JÚNIOR, Ivan; POLÍCIA CIVIL, Góis;
BONTEMPO, Márcio; SILVA, Marcos Roberto; BARROS, Pedro Octavio de Souza.
\textit{Segurança Alimentar e Agronegócio Sustentável no Estado de São Paulo}.
São Paulo: ADESG, 2026.

\vspace{0.5cm}

Trabalho de Conclusão de Curso -- Associação dos Diplomados da Escola Superior
de Guerra, Delegacia do Estado de São Paulo, 2026.

\vspace{0.5cm}

1. Segurança alimentar. 2. Agronegócio sustentável. 3. Planejamento estratégico.
4. São Paulo. I. Associação dos Diplomados da Escola Superior de Guerra.
\end{minipage}}
\end{center}
\end{fichacatalografica}
% ---

% ---
% Folha de aprovação
% Preencha a data, os nomes e as instituições da banca examinadora.
% ---
\begin{folhadeaprovacao}
\begin{center}
\textbf{FOLHA DE APROVAÇÃO}
\end{center}

\noindent
Trabalho de Conclusão de Curso apresentado à Associação dos Diplomados da
Escola Superior de Guerra, Delegacia do Estado de São Paulo, como requisito
para a obtenção do diploma do Curso de Estudos de Política e Estratégia.

\vspace{1cm}

\noindent
Aprovado em: \underline{\hspace{1.5cm}}/\underline{\hspace{1.5cm}}/\underline{\hspace{2cm}}.

\vfill

\assinatura{Prof. Reinaldo Bento de Lima\\Orientador}
\assinatura{Profa. Nídia Caivano\\ADESG, Delegacia do Estado de São Paulo}
\assinatura{Nome do membro da banca\\Instituição}
\end{folhadeaprovacao}
% ---

% ---
% Anverso da folha de rosto:
% ---

% {
% \ABNTEXchapterfont

% \vspace*{\fill}

% Conforme a ABNT NBR 10719:2015, seção 4.2.1.1.1, o anverso da folha de rosto
% deve conter:

% \begin{alineas}
%   \item nome do órgão ou entidade responsável que solicitou ou gerou o
%    relatório; 
%   \item título do projeto, programa ou plano que o relatório está relacionado;
%   \item título do relatório;
%   \item subtítulo, se houver, deve ser precedido de dois pontos, evidenciando a
%    sua subordinação ao título. O relatório em vários volumes deve ter um título
%    geral. Além deste, cada volume pode ter um título específico; 
%   \item número do volume, se houver mais de um, deve constar em cada folha de
%    rosto a especificação do respectivo volume, em algarismo arábico; 
%   \item código de identificação, se houver, recomenda-se que seja formado
%    pela sigla da instituição, indicação da categoria do relatório, data,
%    indicação do assunto e número sequencial do relatório na série; 
%   \item classificação de segurança. Todos os órgãos, privados ou públicos, que
%    desenvolvam pesquisa de interesse nacional de conteúdo sigiloso, devem
%     informar a classificação adequada, conforme a legislação em vigor; 
%   \item nome do autor ou autor-entidade. O título e a qualificação ou a função
%    do autor podem ser incluídos, pois servem para indicar sua autoridade no
%    assunto. Caso a instituição que solicitou o relatório seja a mesma que o
%    gerou, suprime-se o nome da instituição no campo de autoria; 
%   \item local (cidade) da instituição responsável e/ou solicitante; NOTA: No
%    caso de cidades homônimas, recomenda-se o acréscimo da sigla da unidade da
%    federação.
%   \item ano de publicação, de acordo com o calendário universal (gregoriano),
%   deve ser apresentado em algarismos arábicos.
% \end{alineas}

% \vspace*{\fill}
% }

% ---
% Agradecimentos
% ---
\begin{agradecimentos}

\end{agradecimentos}
% ---

% ---
% RESUMO
% ---

% resumo na língua vernácula (obrigatório)
\setlength{\absparsep}{18pt} % ajusta o espaçamento dos parágrafos do resumo
\begin{resumo}
 Segundo a \citeonline[3.1-3.2]{NBR6028:2003}, o resumo deve ressaltar o
 objetivo, o método, os resultados e as conclusões do documento. A ordem e a extensão
 destes itens dependem do tipo de resumo (informativo ou indicativo) e do
 tratamento que cada item recebe no documento original. O resumo deve ser
 precedido da referência do documento, com exceção do resumo inserido no
 próprio documento. (\ldots) As palavras-chave devem figurar logo abaixo do
 resumo, antecedidas da expressão Palavras-chave:, separadas entre si por
 ponto e finalizadas também por ponto.

 \noindent
 \textbf{Palavras-chaves}: latex. abntex. editoração de texto.
\end{resumo}
% ---

% ---
% inserir lista de ilustrações
% ---
\pdfbookmark[0]{\listfigurename}{lof}
\listoffigures*
\cleardoublepage
% ---

% ---
% inserir lista de tabelas
% ---
\pdfbookmark[0]{\listtablename}{lot}
\listoftables*
\cleardoublepage
% ---

% ---
% inserir lista de abreviaturas e siglas
% ---
\begin{siglas}
  \item[ABNT] Associação Brasileira de Normas Técnicas
  \item[abnTeX] ABsurdas Normas para TeX
\end{siglas}
% ---

% ---
% inserir lista de símbolos
% ---
\begin{simbolos}
  \item[$ \Gamma $] Letra grega Gama
  \item[$ \Lambda $] Lambda
  \item[$ \zeta $] Letra grega minúscula zeta
  \item[$ \in $] Pertence
\end{simbolos}
% ---

% ---
% inserir o sumario
% ---
\pdfbookmark[0]{\contentsname}{toc}
\tableofcontents*
\cleardoublepage
% ---


% ----------------------------------------------------------
% ELEMENTOS TEXTUAIS
% ----------------------------------------------------------
\textual

% ----------------------------------------------------------
% Introdução (exemplo de capítulo sem numeração, mas presente no Sumário)
% ----------------------------------------------------------
\chapter*[INTRODUÇÃO]{INTRODUÇÃO}
\addcontentsline{toc}{chapter}{INTRODUÇÃO}

% ---
% Fase de diagnóstico
% ---
\chapter{Fase de diagnóstico}

\section{Análise do ambiente}
\subsection{Antecedentes}
\subsection{Semente de Futuro}
\subsubsection{Fatos Portadores de Futuro}
\subsubsection{Eventos Futuro Preliminares}
\subsubsection{Tendências de Peso}
\subsubsection{Mega-tendências}
\subsubsection{Atores Sociais ou Stakeholders}
\subsection{Ambientes Externo e Interno}
\subsubsection{Análise do Ambiente Externo}
\subsubsection{Análise do Ambiente Interno}
\subsection{Análise da Matriz SWOT Entre-cruzada e Vetores Estratégicos}
\subsubsection{Análise dos Resultados do Entrecruzamento da Matriz SWOT}
\subsection{Análise dos Riscos}
\subsubsection{Análise de Riscos das Ameaças e das Oportunidades}

% \section{Análise do poder}
% \subsection{Planos em Vigor}
% \subsection{Necessidades}
% \subsection{Meios Disponíveis e Potenciais}

% \chapter{Fase política}
% \section{Elaboração de cenários}
% \subsection{Análise Estrutural}
% \subsubsection{Reavaliação das Tendências de Peso}
% \subsubsection{Análise dos Impactos Cruzados}
% \subsubsection{Análise da Estratégia dos Atores}
% \subsubsection{Matriz de Interesse dos Atores}
% \subsection{Elaboração de Cenários Extremos}
% \section{Concepção política}
% \subsection{Objetivos de Estado e Objetivos de Governo}
% \subsection{Fatores Críticos de Sucesso}

% \chapter{Fase estratégica}
% \section{Concepção estratégica}
% \subsection{Levantamento dos Óbices}
% \subsection{Definição das Linhas de Ação Estratégicas}
% \subsubsection{Teste de Adequabilidade}
% \subsubsection{Teste de Exequibilidade}
% \subsubsection{Teste de Aceitabilidade}
% \subsection{Opção Estratégica}
% \subsection{Diretrizes Estratégicas}
% \section{Programação}
% \subsection{Orçamentos}
% \subsection{Planos, Programas e Projetos}

% \chapter{Fase da gestão}
% \section{Execução}
% \subsection{Coordenação}
% \subsection{Implantação}
% \section{Controle}
% \subsection{Acompanhamento}
% \subsection{Avaliação}

% \chapter{Conclusão}

% ----------------------------------------------------------
% ELEMENTOS PÓS-TEXTUAIS
% ----------------------------------------------------------
\postextual

%---------------------------------------------------------------------
% Bibliografia
%---------------------------------------------------------------------
\renewcommand{\bibname}{BIBLIOGRAFIA}
\bibliography{tcc-grupo3-charlie}

%---------------------------------------------------------------------
% Componentes do grupo
%---------------------------------------------------------------------
\chapter{Componentes do grupo}
\label{componentes-grupo}

\textbf{Componentes do grupo}

\bigskip

\small
\begin{tabularx}{\textwidth}{|p{1.7cm}|>{\raggedright\arraybackslash}X|p{3.2cm}|p{3cm}|p{2.2cm}|}
\hline
\textbf{No. de inscrição} & \textbf{Nome} & \textbf{Título profissional} & \textbf{Função no grupo} & \textbf{Assinatura} \\
\hline
06 & Cícero Santos Melo & & Pesquisador & \\
\hline
30 & Ivan Lacava Júnior & & Pesquisador & \\
\hline
16 & Góis Polícia Civil & & Pesquisador & \\
\hline
20 & Márcio Bontempo & & Pesquisador & \\
\hline
23 & Marcos Roberto Silva & & Pesquisador & \\
\hline
26 & Pedro Octavio de Souza Barros & & Pesquisador & \\
\hline
\end{tabularx}

\end{document}
